% !PRINTED 274
\weberpage{274}
%<<CONT
immer zwei Wurzeln hat. Diese beiden Wurzeln sind von einander
verschieden, wenn wir Substitutionen $p^{\mathrm{ten}}$ Grades ausschliessen,
die ja in einer Gruppe, deren Grad nicht durch $p$ theilbar ist, nicht
vorkommen können, und die nach \S.\,68, 1.\ die einzigen sind, für welche
die beiden Wurzeln von (2) zusammenfallen.

Diese beiden Wurzeln nennen wir die Pole von $\Theta$.

Wir untersuchen eine Gruppe $G$, deren Grad $n$ nicht durch $p$
theilbar ist, die ein Theiler von $L_p$ sein soll. Wenn in $G$ die
Substitutionen
\[
\Theta_1,\ \Theta_2,\ \ldots,\ \Theta_{\nu-1},
\]
aber keine anderen vorkommen, die denselben Pol $\alpha$ haben, so
nennen wir $\alpha$ \gesperrt{einen $\nu$-zähligen Pol} (\S.\,52).

Ist nun $S$ irgend eine Substitution der Gruppe $\mathfrak{E}_{2,p}$, so
erhalten wir, wenn wir $L_p$ durch $S^{-1}$ transformiren, eine mit $L_p$
isomorphe Gruppe $S L_p S^{-1}$, und $G$ geht durch dieselbe
Transformation in eine isomorphe Gruppe $S G S^{-1}$ über.

Obwohl die Substitutionen dieser letzteren Gruppe keine reellen
Co\"efficienten haben, so hat doch jede von ihnen zwei Pole, die man
erhält, wenn man die Substitution $S$ auf die Pole von $\Theta$
anwendet. Denn es ist, wenn $\alpha$ ein Pol von $\Theta$ ist:
\[
S\,\Theta\,S^{-1}\,S(\alpha) = S\,\Theta(\alpha) = S(\alpha).
\]
Wir haben also, wie in \S.\,51, 2.:

\gesperrt{1. Die Gruppe $G$ lässt sich so transformiren, dass eine
beliebige ihrer Substitutionen die Pole 0 und $\infty$ erhält, dass
also diese Substitution multiplicativ wird.}

Wir führen nun der Reihe nach die Sätze des \S.\,52 auf, soweit sie
hier in Frage kommen, und werden nur da, wo die veränderten
Voraussetzungen es erfordern, eine Ausführung hinzufügen:

\gesperrt{2. Ist $\alpha$ ein $\nu$-zähliger Pol, so bilden die
Substitutionen 1, $\Theta_1$, $\Theta_2$, . . ., $\Theta_{\nu-1}$ eine
cyklische Gruppe. Ihre Substitutionen können alle (durch
Transformation) in die Form gesetzt werden:}
\weq{3}{\Theta(\xi) = \gamma^{h\frac{p^2-1}{\nu}}\,\xi, \quad h = 0, 1, \ldots, \nu-1 \quad (\S.\,52, 3.).}
